\section{Uniform estimates and widths}
\label{b:width-estimates}

The tower controls discrepancies one relative field at a time. We now
translate this control into estimates for sections: upper bounds on the
dimensions of their initial spaces, together with independent sections
of comparable cost. These two directions are both used in the later
valuation argument. Divisor orders provide the upper bounds; opposite
torus twists and birational systems provide the sections.

Index the tower by \(i=1,\ldots,f\), with
\[
 k(S_i)=k(Y_i)=K_i,\qquad
 Y_i\longrightarrow S_{i+1},\qquad
 S_{f+1}=\mathrm{pt},\qquad K_{f+1}=k=\C.
\]
The first field is the residual field in \eqref{eq:A9}.
After a subsequence all dimensions and ranks are constant. Put
\[
 p_i=\trdeg_{K_{i+1}}K_i,\qquad
 a_i=s_iC_i\asymp s_iA_i^++\Lambda_i,\qquad
 s_f=1,\quad s_i\longrightarrow0\quad(i<f).
\]
Indeed the tower gives, with the constants chosen at level \(i\),
\[
 a_i=(1-\gamma_i)s_iA_i^+
           +b\lambda_iQ_i+s_i(a+x_i)H_i,
\]
omitting the \(Q_1\)-term. All terms have nonnegative values.
Moreover
\[
 \lambda_{i+1}=1/x_i\asymp s_i,\qquad s_{i+1}\gtrsim s_i.
\]
Constants in this section may depend on the previously fixed
sequence bounds and dimensions, but are uniform in the functions,
divisors, weights and cutoffs subsequently tested.

The b-nef Calabi--Yau moduli of \(a_i\) contain simultaneous
positive reserves of \(\mathbf H_i=\overline{H_i}\) and, for \(i>1\),
of \(s_{i-1}\mathbf M_{Q_i}\). The first coefficient is
\(s_i(x_i-d_i')\gtrsim1\); the second follows from
\(s_i\gamma_i=b\lambda_i\).
For an invariant prime \(E\) on \(\Spec R\), write
\(E|_{K_i}=j_iP_i\) when nontrivial, with \(P_i\) normalized. Then
\begin{equation}\label{b:scaled-witness}
                         j_ia_i(P_i)\le s_i+C_0h_E.
\end{equation}
To see the exact constant, substitute \eqref{eq:B2} in the displayed
formula for \(a_i\). The constant contributions are
\((1-\gamma_i)s_i+b\lambda_i=s_i\); the remaining terms are bounded
multiples of \(h_E\), using \(j_iH_i(P_i)\le h_E\).
The unscaled $C_i$ have uniform positive horizontal discrepancy bounds.
The scaled functions $a_i=s_iC_i$ measure the smaller divisor orders
that correspond to width $1/s_i$; their positive lower bounds are not
uniform at collapsing levels. All restrictions and discrepancy functions
are evaluated homogeneously. At a trivial restriction their value is defined to
be zero.

\subsection{Horizontal orders and errors of nef traces}

Let \(\mathcal F_i\) be the arithmetic generic fiber of
\(Y_i\to S_{i+1}\), over \(K_{i+1}\); for \(i=f\) this is \(Y_f\).
The models \(Y_i\) are \(\Q\)-factorial. A horizontal prime means
a prime-divisor valuation on \(K_i\) trivial on \(K_{i+1}\).
Horizontally \(C_i=a_{K_i}+\mathbf M_{C_i}\) has effective uniformly
positive-lc generalized Calabi--Yau data with big nef part.
For \(i<f\), use the transformed threshold discrepancy
\(C_i-x_i\pi^*H_{i+1}\); for \(i=f\), use the last-step data.
The trace boundary \(B_{0,i}''\) of
\(A_i^+=a_{K_i}+\mathbf M_{0,i}\) is effective and dominates the
trace \(H_i''\) of \(\mathbf H_i\).

The normal projective arithmetic fibers are geometrically integral
and normal, by characteristic zero and the contraction property.
Their arithmetic Weil divisors are rationally Cartier, by taking
closures on \(Y_i\). We do not require geometric
\(\Q\)-factoriality. Divide the chosen volume form by a base volume
form generating canonical on a smooth base open. At horizontal
generic points, smoothness of the generic fiber identifies the
resulting relative canonical coefficients, including on resolving
models. The dimension formula identifies normalized horizontal
tests with prime-divisor tests of the relative field.

Such a test extends after algebraic closure of \(K_{i+1}\) with
unchanged normalization on arithmetic functions: take a component
of its extension on a smooth birational chart, whose generic
ramification index is one by separability. Conversely snc
resolution transfers the arithmetic discrepancy inequalities to all
geometric tests. In particular an snc boundary with coefficients at
most \(1-\delta\), \(0<\delta\le1\), has discrepancy at least
\(\delta\) on every normalized prime: subtract the reduced-support
discrepancy and use integrality of the remaining orders.

Lemma~\ref{b:ordinary-approximation}, with an ample reserve, makes
the geometric fibers klt of Fano type and puts them in a bounded
family by \cite[Theorem 1.1 and Corollary 1.4]{BirkarBAB}.
Use a bounded very ample Cartier polarization \(D\) to measure
degrees. These bounded-family assertions over \(\C\) apply here:
the algebraic closure of a finitely generated function field over
\(\C\) is individually isomorphic to \(\C\).
The ordinary effective approximation \(\Delta\) can be chosen with
\(K+\Delta\sim_\Q0\), a uniform positive-lc bound, and ordinary
discrepancies no larger than \(C_i\). This approximation may first
be made on the whole threshold model and then restricted
geometrically. Dimension-zero fibers need none of the following
divisor estimates.

\begin{lemma}\label{b:horizontal-estimates}
For every effective rational \(\Q\)-Cartier divisor \(L'\) on
\(\mathcal F_i\) and every horizontal test \(P\),
\begin{equation}\label{eq:C1}
              P(L')\le C_0\deg(L')C_i(P).
\end{equation}
If \(\mathbf U\) is rational b-Cartier and b-nef, with trace
\(U_{\rm tr}\) on \(Y_i\), then
\begin{equation}\label{eq:C2}
 0\le P(\overline{U_{\rm tr}}-\mathbf U)
                  \le C_0\deg(U_{\rm tr})C_i(P).
\end{equation}
\end{lemma}

\begin{proof}
Put \(p=p_i>0\). The bare geometric fiber is klt, so the cone
theorem and its length bound \cite[Theorem~1.1(5)]{FujinoFundamental}
give \(K+2pD\) nef. Thus \(A=(2p+1)D\) satisfies
\(A-\Delta\sim_\Q A+K\) ample. Its degree is bounded, and the
polarized threshold theorem \cite[Theorem 1.8]{BirkarBAB}
gives a uniform positive lower threshold for the
\(\R\)-linear system of \(A\) against the ordinary positive-lc pair.
The same cone inequality bounds \(\deg(-K)\).

Every geometric prime of degree \(e\) lies in a Cartier section of
\(|qD|\) for some \(1\le q\le e\). Indeed a general basepoint-free
projection to \(\PP^p\) is finite, since its pulled-back hyperplane
is ample. The image of the prime is a hypersurface of degree at
most \(e\), and its pullback contains the prime. Consequently a
nonzero \(L'\) is dominated geometrically by a rational weighted
sum of such sections, of rational class \(q'D\) with
\(q'\le\deg L'\). The difference is effective rationally Cartier,
even when the individual geometric primes are not rationally
Cartier. Scale this sum to the system of \(A\), apply the threshold
bound, and compare the ordinary discrepancy with \(C_i\).
Extension of the arithmetic test gives \eqref{eq:C1}.

Negativity makes \(\overline{U_{\rm tr}}-\mathbf U\) b-effective.
On a high projective descent add a positive rational ample amount
tending to zero. The error of the ample addition is again
nonnegative, so this can only increase the error being bounded.
The perturbed nef class is ample and has, after a rational
principal adjustment, an effective representative upstairs.
That adjustment changes neither error. Its nonnegative orders
bound the new error by the pullback order of its effective trace.
Apply \eqref{eq:C1} and let the added amount tend to zero.
The traces in question are arithmetically rationally Cartier
and geometrically pseudo-effective; their degrees are
nonnegative. This proves \eqref{eq:C2}.
\end{proof}

\subsection{A pure upper estimate}

\begin{proposition}\label{b:pure-upper}
For \(D_1\ge0\), \(D_1\sim_\Q H_1\) on \(S_1\), and every divisorial
test \(P_1\), with restriction multiplicities included in
\(P_i=P_1|_{K_i}\), one has
\begin{equation}\label{eq:C3}
                         D_1(P_1)\le C_0\sum_i a_i(P_i).
\end{equation}
For tests trivial on \(K_f\), this implies
\begin{equation}\label{eq:C4}
 D_1(P_1)\le C_0\bigl(H_1(P_1)+s_{\exp}A_1^+(P_1)\bigr),
 \qquad s_{\exp}:=\max_{i<f}s_i\longrightarrow0,
\end{equation}
where an empty maximum is zero.
\end{proposition}

\begin{proof}
Horizontally on \(Y_1\), \(\deg D_{1,\rm tr}\lesssim s_1\).
For a nonlast level, \(x_1H_1''=J''\le B_{0,1}''\);
at the last level use \(H_1''\le B_{0,1}''\).
The Calabi--Yau trace identity and nonnegative moduli degree bound
\(\deg B_{0,1}''\) by \(\deg(-K)\).
The b-nef divisor \(D_1\) has orders at most those of its effective
trace closure. Hence \eqref{eq:C1} gives
\(D_1(P)\le C_0a_1(P)\) on horizontal tests.
Subtract \(\tau_1D_1\) from \(a_1\) for a small uniform positive
rational \(\tau_1\). The moduli stay nef and big by the
\(\mathbf H_1\)-reserve, and the horizontal data stay generalized
klt and effective.

Inductively define
\[
 D_{i+1}=\pi_\#a_i-\pi_\#(a_i-\tau_iD_i).
\]
The minimum formula makes it nonnegative everywhere.
Proposition~\ref{b:push-presentation} writes it as
\(\mathbf U-\mathbf U'\) with both terms rational b-Cartier
and b-nef. Their canonical coefficients agree, and
\[
 \mathbf U=s_i\mathbf M_{Q_{i+1}},
 \qquad
 \pi_\#C_i=Q_{i+1}+x_iH_{i+1}.
\]
The next level has a positive reserve of this \(\mathbf U\).
Its effective scaled horizontal Calabi--Yau identity therefore
bounds \(\deg U_{\rm tr}\le C_0s_{i+1}\).
Nonnegativity of trace degrees and of the trace of \(D_{i+1}\)
gives the same bound for \(U'_{\rm tr}\) and \(D_{i+1,\rm tr}\).
Apply \eqref{eq:C1} to the latter trace and \eqref{eq:C2} to
the two trace errors. Thus
\[
                   0\le D_{i+1}(P)\le C_0a_{i+1}(P)
\]
horizontally. For a uniform sufficiently small \(\tau_{i+1}>0\),
subtracting \(\tau_{i+1}D_{i+1}\) preserves nef big moduli:
it subtracts \(\tau_{i+1}\mathbf U\) from its reserve and adds
\(\tau_{i+1}\mathbf U'\). It also preserves horizontal klt
and trace effectivity. A level without horizontal tests needs
no degree estimate.

For a test continuing nontrivially to \(P_{i+1}\),
\[
 (a_i-\tau_iD_i)(P_i)
 \ge \pi_\#(a_i-\tau_iD_i)(P_{i+1})
 \ge -D_{i+1}(P_{i+1}),
\]
because \(a_i\ge0\). Rearranging and using the horizontal estimate
at the last nontrivial restriction proves \eqref{eq:C3}.
For a test trivial on \(K_f\), only \(i<f\) contribute.
The tower gives
\[
 A_{i+1}^+(P_{i+1})\le A_i^+(P_i),\qquad
 \Lambda_{i+1}(P_{i+1})
 \le\lambda_{i+1}C_i(P_i)\lesssim a_i(P_i).
\]
Recursively use \(a_i\asymp s_iA_i^++\Lambda_i\) and
\(\Lambda_1=H_1\); the tower length is bounded.
This gives \eqref{eq:C4}.
\end{proof}

For tests trivial on the terminal field $K_f$, \eqref{eq:C4} involves
only the scales that tend to zero. The uniformly bounded residual
directions have disappeared from that estimate.

\subsection{Integral classes for opposite characters}

Recall \(M_T=\Z\chi\oplus M_0\), with \(M_0=\ker D\), and the
normalized dual polytope \(Q_T\). On \(M_{0,\R}\) set
\[
                       \|\beta\|=\max_{\rho\in Q_T}|\rho(\beta)|.
\]
This is a norm by the occurring-cone and interiority properties.

For an occurring weight \(\mu=m\chi+\beta\) of positive cost
\(B'=\widetilde v(\mu)\), the comparison on \(Q_T\) following
\eqref{eq:A12} gives
\[
                         0\le m\lesssim B',\qquad
                         \|\beta\|\lesssim B'.
\]
Indeed \(D\in Q_T\) and \(0\le\rho(\mu)\le CB'\) for all
\(\rho\in Q_T\).

To see how the pure estimate will control a general weight space, write
$R_\mu=z^\mu U$ with coefficient space $U\subset K_1$.
Finite generation supplies an opposite section of weight
$d(m'\chi-\beta)$ for a suitable $m'>\|\beta\|$ and some positive
clearing integer $d$. If its coefficient is $g'$, then for every
$g\in U\setminus\{0\}$ the product
$z^{d(m+m')\chi}g^dg'$ is a pure section. Both its degree and the
exponent of $g$ have been multiplied by $d$, so this factor cancels
when \eqref{eq:C3} is used to bound the order of $g$. We carry out the
resulting coefficient count below.

Actual probes require more: their degrees and costs must be bounded by
a uniform multiple of $\max\{1,\|\beta\|\}$, so their clearing factor
must be bounded as well.
The following integral class data make that possible on the bounded
fibres. They will also control the divisor classes in the marked families
of Section~\ref{b:bounded-axes}.

\begin{lemma}[Integral class data]\label{b:class-lattices}
For projective normal klt varieties of Fano type occurring in a
bounded family with bounded polarizations, the groups
\(\Cl/\mathrm{torsion}\) can be identified with Euclidean lattices
of finitely many types, so that
\begin{equation}\label{b:class-norm}
                   \|[G]\|_{\rm cl}\le C_0\deg(G)
\end{equation}
for effective rational Weil divisors and for pseudo-effective
rationally Cartier rational classes. The orders of \(\Cl\)-torsion
are bounded. For the \(\Q\)-factorial members, a single bounded
positive multiple of every integral Weil divisor is Cartier.
A numerically trivial rational Cartier divisor is rationally
linearly trivial.
\end{lemma}

\begin{proof}
The bounded family need only contain the varieties under
consideration; its other members need not satisfy positivity.
After finite stratification and finite base extensions, take
simultaneous projective smooth resolutions \(b_t:W_t\to X_t\)
with labeled exceptional prime divisors. To construct them,
resolve the geometric generic fiber
\cite[Theorems~1.0.1--1.0.3]{WlodarczykResolution}, spread after a
finite extension, shrink, and apply noetherian induction to the
complement. Keep an isomorphism on dense opens and spread the
components of the complements. Stabilizing their fiber and image
dimensions identifies all exceptional prime families and excludes
additional exceptional divisors. Smoothness of \(W_t\), normality
and geometric integrality of the fibers in use, and geometric
integrality of the prime families hold after the corresponding
open restrictions. Fixed labeled prime or Cartier data can be
included in this process.

On the selected fibers, Kawamata--Viehweg vanishing for a klt log
Fano boundary \cite[Theorem~3.2]{FujinoFundamental} and rationality
of klt singularities \cite[Theorem~0.1]{FujinoRational} give
\(H^1(\OO)=H^2(\OO)=0\), both on \(X_t\) and on \(W_t\).
The exponential sequence and GAGA
\cite[Section~12, Theorems~1--3; Section~20, Proposition~18]{SerreGAGA}
therefore identify \(\Pic\) with integral \(H^2\) on both. Strict transform and pushforward give the diagram
\[
 H^2(X_t,\Z)\ \xrightarrow{\,b_t^*\,}\ H^2(W_t,\Z)
 \ \longrightarrow\
 H^2(W_t,\Z)/\langle[E_1],\ldots,[E_q]\rangle=\Cl(X_t).
\]
The image of the first group in the quotient is \(\Pic(X_t)\).

These are locally constant integral group data, including torsion,
after stratification: proper topological constructibility gives
the integral cohomology local systems and functorial pullback,
while relative line bundles give the labeled Chern classes.
Here is the finite stratification argument, including compatibility
with the markings. For either projective family \(\mathcal V\to T\),
let \(Z_1,\ldots,Z_s\subset\mathcal V\) be its finitely many closed
marked subsets. Apply \cite[Theorem~3.2.3]{DeCataldoMiglioriniHodge}
to the finite map
\[
 \mathcal V\sqcup\bigsqcup_{i=1}^s Z_i\longrightarrow\mathcal V
\]
given by the identity and the inclusions. Split the source strata by
the open-and-closed summands and their connected components. For each
connected target stratum \(S\), the subset \(Z_i\cap S\) is closed.
It is also open, since it is a union of images of submersive source
strata. Hence it is either empty or all of \(S\), so the finite target
Whitney stratification is compatible with every marking.

Work over each irreducible base piece and restrict to a smooth dense
open. Remove the closures of the
images of all nondominating strata and of the critical loci of the
dominating strata. These closures are proper by generic smoothness.
The remaining strata map submersively to the base, and the whole
family remains proper. Embed the family relatively in a smooth
projective ambient space. Thom's first isotopy lemma
\cite[Proposition~11.1]{MatherTopological} gives local topological
triviality. Intersect the opens for the finitely many families in use.
Noetherian induction treats the remaining base pieces by reapplying
this construction to their reduced base-changed families and closed
markings; no Whitney assertion for arbitrary closed restrictions is
needed.

In trivializations over a small contractible base open, the fiber
maps vary continuously and hence induce constant cohomology maps.
No simultaneous trivialization of the resolution morphism, or trivial
monodromy, is required. Monodromy preserves the isomorphism type of
this group diagram.

There are consequently finitely many group-diagram types.
Their cokernels bound the torsion of \(\Cl\); when
\(\Cl/\Pic\) is finite, its exponent is bounded as well.
The latter finiteness is exactly rational factoriality here,
so it gives the asserted Cartier multiple.

For the norm assertion in dimension \(p\ge2\), fix a relative very
ample \(D_W\) on each resolved stratum. After further finite
extension and shrinking, choose relative line bundles
\(T_1,\ldots,T_r\) spanning \(H^2(W_t,\Q)\), with a common integer
\(l>0\) such that \(lD_W\pm T_j\) are ample. We justify this choice,
which contains more information than the abstract local systems.
Shrink where the coherent vanishings hold and the Betti ranks are
constant. Descend the finite family data to a countable
algebraically closed subfield of \(\C\), and realize its geometric
generic fiber as a complex fiber generic over that smaller field.
The relevant algebraic closures and \(\C\) are isomorphic over the
embedded defining parameter field, since their remaining
transcendence degrees have the same infinite cardinality.
On that fiber \(\Pic=H^2(\Z)\). Hard Lefschetz and the
Hodge--Riemann bilinear relations
\cite[Theorem~3.1.2(a)--(c)]{DeCataldoMiglioriniHodge} make the
pairing \(TT'D_W^{p-2}\) nondegenerate on \(H^2(\Q)\).
Indeed \(H^2=P^2\oplus D_WH^0\) is an orthogonal decomposition:
the primitive summand is nondegenerate by polarization, and the
ample summand has nonzero self-pairing \(D_W^p\).
Choose spanning line bundles there, spread them after a finite
extension, and retain their nonsingular intersection matrix and
ample comparisons on an open. Betti-rank constancy proves they
continue to span. Noetherian induction treats the other strata.

Use the coordinates
\[
                  \bigl(G_W\cdot T_j\cdot D_W^{p-2}\bigr)_j.
\]
For an effective divisor their absolute values are at most
\(l\deg_{D_W}(G_W)\), by the ample comparisons.
The integral free lattice embeds in \(\Z^r\) and contains the
fixed full-rank subgroup generated by the \(T_j\)-coordinates.
Thus only finitely many intermediate lattices occur. The
exceptional classes have fixed coordinates too. Orthogonal
projection modulo their real span gives lattices of finitely
many types for the free class groups downstairs.

If \(G_W\) is the strict transform of an effective prime \(G\),
the containing hypersurface used in the proof of \eqref{eq:C1}
pulls back to an effective Cartier divisor dominating \(G_W\).
The fixed family intersection degrees give
\(\deg_{D_W}(G_W)\le C\deg(G)\).
Linearity and the triangle inequality prove \eqref{b:class-norm}
for effective rational Weil divisors. In dimension one use
ordinary degree and the Picard description above; dimension
zero is vacuous.

A numerically trivial rational Cartier divisor pulls back to
a numerically trivial class upstairs. The same nondegenerate
pairing and \(\Pic(W_t)=H^2(W_t,\Z)\) make its rational class zero.
Pushforward proves rational linear triviality downstairs.
Finally a pseudo-effective rational Cartier class becomes big
after adding any small positive rational ample amount, and a
big rational Cartier class is rational-linearly effective.
Apply the effective bound and pass to the limit.
\end{proof}

\begin{lemma}[Arithmetic principal witnesses]\label{b:principal-descent}
An integral Weil divisor defined on a proper geometrically integral
normal arithmetic fiber which becomes principal geometrically is
already principal over the arithmetic field. No extension-degree
factor is needed. The same holds for a rational divisor after
clearing its specified denominator.
\end{lemma}

\begin{proof}
A witnessing function is defined over a finite normal extension
\(L/K\). For \(\sigma\in\operatorname{Gal}(L/K)\), its conjugate
quotient has zero divisor, hence is an invertible global regular
function on the proper normal fiber. Geometric integrality makes
it a constant in \(L^*\). The resulting scalar cocycle is trivial
by Hilbert 90. Rescale the function by the corresponding scalar;
its divisor is unchanged and the function is Galois invariant.
\end{proof}

\subsection{Saturation of opposite twists}

For a character $\beta$ with $\|\beta\|\le L$ and $L\ge1$, we now use
the integral class data to replace arbitrary monoid clearing by a bounded
factor. On each fibre for which $Ls_i\to0$, discreteness
of divisor classes forces the obstruction to be a multiple of the pivot
class. A rational principal correction removes it, and the discrepancy
bounds allow us to pass the corrected inequalities to the next field.
At the first scale with $Ls_i\gtrsim1$, a section construction and
integrality close the argument.

\begin{proposition}\label{b:twist-saturation}
If \(\beta\in M_0\), \(\|\beta\|\le L\), and \(L\ge1\), then
there are \(m\in n\Z_{\ge0}\) and a bounded positive integer
\(d_0\), with \(m\lesssim L\), such that
\begin{equation}\label{eq:C5}
                R_{m\chi+d_0\beta}\ne0,\qquad
                R_{m\chi-d_0\beta}\ne0 .
\end{equation}
\end{proposition}

\begin{proof}
We may test the assertion on a contrary sequence, allowing $\beta$ and
$L$ to vary with its member, and take subsequences throughout.
The first clearing factor $d$ below may be arbitrary; only the eventual
factor $d_0$ is required to be uniformly bounded.
At each individual member, finite
generation of the weight monoid gives the assertion with
member-dependent constants: rational cone points enter the monoid
after a common multiple, by decomposition into simplicial
generator cones. For \(m_0\in n\Z\) a little larger than \(L\),
both \(m_0\chi\pm\beta\) lie in the cone by its normalized dual
inequalities. Thus \(m_0\lesssim L\), and actual opposite sections
exist after a possibly arbitrary common clearing factor \(d>0\):
\[
                 \bar h^{dm_0/n}z^{\pm d\beta}r_\pm,
                     \qquad r_\pm\in K_1^* .
\]
We can take late members of a contrary sequence and round \(L\)
up to an integer. Start with
\[
         u_\pm=-\Div(r_\pm)/d,\qquad
         \ell_E=E(z^\beta),\qquad \mu=0.
\]
The arbitrary \(d\) occurs in \(u_\pm\); the initial
\(\ell_E\) are integral.

After subsequences, stop at the first stage with \(Ls_i\gtrsim1\);
at all preceding stages \(Ls_i\to0\). A stop exists because
\(s_f=1\). At entry to each reached stage we maintain
\(\mu\ge0\), \(\mu=o(1)\), \(\mu\lesssim L\lambda_i\), and
\begin{equation}\label{eq:C6}
\begin{aligned}
 u_\pm&=\mu a_{K_i}+\mathbf U_\pm,
 &\mathbf U_\pm&\text{ rational b-Cartier and b-nef},\\
 u_++u_-&\le C_0L\Lambda_i,\\
 j_iu_\pm(P_i)&\le\pm\ell_E+C_0Lh_E+\mu,
 &j_iP_i&=E|_{K_i}.
\end{aligned}
\end{equation}
The first inequality is on every prime and is homogeneous.
The last line has left side zero at a trivial restriction.
We allow \(\ell_E\) to change by orders of fractional rational
functions from tower fields, with bounded clearing denominators.
Such fractional functions are notation for rational principal
b-divisors, not field extensions.
Initially the product of the two sections and \eqref{eq:A9} give
\[
                  \Div(r_+r_-)+2dm_0H_1\ge0,
\]
which is the sum bound; their regularity gives the last line.

\smallskip\noindent
\emph{Removing the class obstruction before the stopping scale.}
Consider a stage before the stop, so $Ls_i\to0$. On its arithmetic fiber every strict
prime other than the pivot \(G_*\) comes from a strict prime of
\(S_i\) with \(H_i(P)=0\): the pivot alone among the positive
support survives horizontally, and every extracted prime had
positive \(H_i\). Choose a witness \(j_iP=E|_{K_i}\) from
\eqref{eq:B2} for each such prime, with bounded positive \(j_i\)
and \(h_E=0\), and set
\[
 c_P=\ell_E/j_i,\qquad
 \Delta_{\rm col}=\sum_{P\ne G_*}c_PP.
\]
This arithmetic divisor has finite support and bounded denominator.
Indeed \(\ell_E\) is initially an integral character order, and
subsequently only the orders of previously bounded fractional
gauges are added. There are dimension-many stages and bounded
restriction integers. The original arbitrary \(d\) never enters
\(\Delta_{\rm col}\). A rational function has nonzero orders at
only finitely many strict primes of \(\Spec R\), proving the
support assertion.

Write \(B_0''=B_{0,i}''\), \(\mathbf M_0=\mathbf M_{0,i}\), and
\(c_\pm^*=u_\pm(G_*)\). At nonpivot strict primes,
\[
 (\mathbf U_\pm-\mu\mathbf M_0)(P)
       =u_\pm(P)-\mu A_i^+(P)
       \le\pm c_P+\mu B_0''(P).
\]
At the pivot \(A_i^+(G_*)=0\), while
\(c_+^*+c_-^*\le C_0Ls_i\). Hence the divisors
\[
 W_\pm=\pm\Delta_{\rm col}+c_\pm^*G_*+\mu B_0''
                         -(U_\pm-\mu M_0)_{\rm tr}
\]
are effective on the arithmetic fiber. Their degrees and those
of the pseudo-effective traces \(U_{\pm,\rm tr}\) are \(O(Ls_i)\):
add the two identities to obtain
\[
 W_++W_-+(U_++U_-)_{\rm tr}
  =(c_+^*+c_-^*)G_*+2\mu(B_0''+M_{0,\rm tr}).
\]
All terms on the left have nonnegative degree.
The degrees of \(B_0''\), \(M_{0,\rm tr}\) and \(G_*\) are bounded
by Calabi--Yau effectivity; the pivot has coefficient one in
\(B_0''\). Also \(\mu\lesssim Ls_i\).

On the geometric fiber, Lemma~\ref{b:class-lattices} applied to
these identities gives
\[
                \|\pm[\Delta_{\rm col}]+c_\pm^*[G_*]\|_{\rm cl}
                         \le C_1Ls_i .
\]
The nonzero integral vector \([G_*]\) has bounded length; it is
nonzero already by its positive degree. Bounded-denominator
lattice vectors projected modulo its line are uniformly discrete,
because there are finitely many lattice types and finitely many
possible bounded pivot vectors. As \(Ls_i\to0\), we obtain the
exact equality
\[
          [\Delta_{\rm col}]=q[G_*],\qquad
                 |c_\pm^*\pm q|\lesssim Ls_i .
\]
The denominator of \(q\) is bounded. In fact if
\([G_*]=h g\) with \(g\) primitive, bounded length bounds \(h\);
if \(Q[\Delta_{\rm col}]\) is integral, then \(Qqh\in\Z\).
After also killing the bounded torsion, a bounded positive
multiple of \(-\Delta_{\rm col}+qG_*\) is principal geometrically.
Lemma~\ref{b:principal-descent} descends the witness without an
extension-degree multiplier. Let \(\phi\) be a fractional function
in \(K_i\), with bounded clearing denominator, having this divisor
on the fiber.

Replace \(u_\pm\) by \(u_\pm\pm\Div(\phi)\), and
\(\ell_E\) by \(\ell_E+E(\phi)\). These principal changes preserve
\eqref{eq:C6}. They give \(u_\pm(P)\le\mu\) on strict nonpivot
primes and \(u_\pm(G_*)\le C_1Ls_i\). The new signed trace of
\(\mathbf U_\pm-\mu\mathbf M_0\) is
\[
       (c_\pm^*\pm q)G_*+\mu B_0''-W_\pm.
\]
Its positive and negative parts have degrees \(O(Ls_i)\), and
the trace degrees of \(\mathbf U_\pm\) are unchanged. All these
arithmetic traces and parts are rationally Cartier.
Apply \eqref{eq:C1} to the signed trace and \eqref{eq:C2} to
the errors of \(\mathbf U_\pm\) and \(\mathbf M_0\), in the identity
\[
 u_\pm=\mu A_i^++\mathbf U_\pm-\mu\mathbf M_0.
\]
Since \(A_i^+\ge0\), this proves, on every horizontal test,
\[
                              u_\pm\ge-C'La_i.
\]

\smallskip\noindent
\emph{Pushing the corrected inequalities.}
The horizontal lower bound just proved will give klt discrepancies after
adding a large multiple of $La_i$. The strict-prime upper bounds give
effective horizontal trace. An additional multiple of $Ls_iA_i^+$
supplies the needed allowance at the pivot.
For sufficiently large fixed positive rational \(K',T'\), push
\[
                    u_\pm+K'La_i+T'Ls_iA_i^+
\]
to the next field and use these pushes as the new \(u_\pm\).
Their canonical coefficient is
\(\mu_{\rm next}=\mu+(K'+T')Ls_i>0\).
They have big nef moduli and horizontal generalized klt
discrepancies by the preceding lower bound.
They are effective in horizontal trace as well:
on strict nonpivot primes use \(u_\pm\le\mu\) and
\(C_i,A_i^+\le1\); at the pivot use \(u_\pm\le C_1Ls_i\),
\(A_i^+=0\), and the extra slack \(T'Ls_i\).
Proposition~\ref{b:push-presentation} therefore applies.
We have \(\mu_{\rm next}=o(1)\) and
\(\mu_{\rm next}\lesssim L\lambda_{i+1}\).

For the sum bound, evaluate both signs on one common lift
minimizing \(C_i\). Use
\(\Lambda_i+s_iA_i^+\lesssim a_i\) and
\(s_i\pi_\#C_i\lesssim\Lambda_{i+1}\).
For the witness bound, if the projection remains nontrivial,
evaluate on the old \(j_iP_i\). By \eqref{b:scaled-witness}
and \eqref{eq:B2}, the newly added constants are exactly
\((K'+T')Ls_i\), while the other terms are \(O(Lh_E)\).
They are therefore precisely absorbed in \(\mu_{\rm next}\).
If the projection becomes trivial, the expression just pushed
has positive value at the old horizontal test, so the same upper
bound holds with left side zero. An already trivial projection
continues to satisfy the bound. This proves the induction.

\smallskip\noindent
\emph{Completing the canonical coefficient and obtaining regular sections.}
At the stopping stage \(Ls_i\gtrsim1\), let \(d_0\) be a bounded
positive integer clearing all accumulated gauges. There are only
boundedly many stages, and eventually \(d_0\mu<1\). The positive
coefficient $1-d_0\mu$ will complete the canonical coefficient to one;
the same equality will give the one-unit allowance in the final integral
regularity test. For each sign
we claim that some \(f_\pm\in K_i^*\) satisfies, at every prime,
\begin{equation}\label{eq:C6a}
 P(f_\pm)+d_0u_\pm(P)+(1-d_0\mu)A_i^+(P)+C_2LH_i(P)>0 .
\end{equation}
If \(S_i\) is a point, take \(f_\pm=1\). Otherwise the middle
terms have the presentation
\[
 a_{K_i}+\mathbf N,\qquad
 \mathbf N=d_0\mathbf U_\pm+(1-d_0\mu)\mathbf M_{0,i},
\]
with b-nef \(\mathbf N\).

On a high smooth descent, \eqref{eq:B3} provides a big free
integral Cartier divisor \(D'\le kH_i\) in actual orders,
with \(k\lesssim1/s_i\lesssim L\).
Indeed choose a list \(1,f_1,\ldots,f_t\) with poles bounded by
\(kH_i\) giving the birational system, resolve its projective map,
and pull back its first coordinate hyperplane. Freeness makes
the order of this divisor \(-\min(0,P(f_1),\ldots,P(f_t))\),
which is at most \(kH_i(P)\).
Resolve the fractional support of \(N=\mathbf N_{\rm tr}\) too,
and use the canonical divisor \(K\) of the same chosen form.
For \(j=1,\ldots,\dim S_i+1\), Kawamata--Viehweg vanishing
\cite[Theorem~3.1]{FujinoFundamental} applies to
\[
                         K+\lceil N\rceil+jD',
\]
with snc klt boundary \(\lceil N\rceil-N\) and nef big adjoint
difference \(N+jD'\). Riemann--Roch makes its Euler
characteristic a polynomial in \(j\) of degree \(\dim S_i\)
with positive leading coefficient. It cannot vanish at all
these integers; vanishing identifies its values with section
dimensions. Choose a nonzero section at one of them.
The associated function \(f_\pm\) makes the left side of
\eqref{eq:C6a}, with \(jD'\) in place of \(C_2LH_i\),
the sum of the positive discrepancy of the snc pair and the
order of the section's effective divisor. This proves the claim.

For a nontrivial \(j_iP_i=E|_{K_i}\), multiply
\eqref{eq:C6a} by \(j_i\) and substitute \eqref{eq:C6},
\(j_iA_i^+(P_i)\le1-h_E\), and \(j_iH_i(P_i)\le h_E\).
The constant allowance is exactly
\[
                         d_0\mu+(1-d_0\mu)=1.
\]
Consequently, for \(m\in n\Z_{\ge0}\) sufficiently large with
\(m\lesssim L\), both signs satisfy
\[
                         E(f_\pm)\pm d_0\ell_E+mh_E>-1.
\]
At trivial restrictions this follows directly from
\eqref{eq:C6} and \(d_0\mu<1\).
Multiply \(f_\pm\) by \(z^{\pm d_0\beta}\), by the corresponding
cleared gauge functions, and by \(\bar h^{m/n}\).
The displayed orders are integers, hence nonnegative.
These are nonzero eigenfunctions. Any polar divisor of an
eigenfunction is torus invariant, so the invariant-prime tests
and normality make them regular. Their weights are
\(m\chi\pm d_0\beta\), proving \eqref{eq:C5} along every
alleged contrary subsequence and therefore uniformly.
\end{proof}

\subsection{Axes, coefficient spaces, and actual probes}

We now combine the residual tower and the torus directions into one list
of widths. We first complete the coefficient count described above and
sum it over the torus characters at each fixed physical degree. That
upper bound uses only arbitrary monoid clearing. The bounded clearing
just proved supplies the matching actual probes and their independent
initials.

Choose a lattice basis \(\beta_1,\ldots,\beta_{\rank T-1}\)
of \(M_0\) and positive widths \(L_j\) so that
\begin{equation}\label{b:torus-axis-norm}
                    \Big\|\sum k_j\beta_j\Big\|
                            \asymp\sum_j|k_j|L_j.
\end{equation}
Here is a uniform construction in bounded dimension.
Replace the norm by an equivalent Euclidean norm using an
ellipsoid, then choose successive shortest nonzero vectors
in the orthogonal projected lattices. Each is primitive, so
the next projection is again a lattice. Choose size-reduced
lifts, subtracting earlier vectors in reverse order until
their earlier Gram--Schmidt coefficients have absolute
value at most \(1/2\). Let \(L_j\) be their Gram--Schmidt
lengths. A size-reduced one-step lift of the next projected shortest vector has
component along the preceding shortest vector at most half its length.
The lift cannot be shorter than that preceding shortest vector, so
\[
              L_j^2\le L_{j+1}^2+(L_j/2)^2.
\]
Thus the next projected length is at least $\sqrt3/2$ times the preceding
shortest length. Hence the normalized basis columns in
orthonormal Gram--Schmidt coordinates form a triangular
matrix with diagonal one and bounded entries; its inverse
is bounded in fixed dimension. This proves
\eqref{b:torus-axis-norm}.

Adjoin \(p_i\) residual axes of width \(1/s_i\) for each level.
Denote all widths by \(W_j\). There are
\[
                    d=(\rank T-1)+\trdeg_kK_1
\]
axes in total.
\begin{proposition}\label{b:coefficient-dimension}
For every occurring weight $\mu=m\chi+\beta$ of positive cost
$B'=\widetilde v(\mu)$,
\begin{equation}\label{eq:C7}
                     \dim R_\mu\lesssim
                           \prod_i(1+B's_i)^{p_i}.
\end{equation}
At each fixed physical degree \(m\), the dimension through
cost \(S'\ge0\), including constants if present, is at most
\begin{equation}\label{eq:C8}
                           C_3\prod_j(1+S'/W_j).
\end{equation}
For \(m\in n\Z_{\ge0}\), the same expression bounds the
truncated associated graded dimension of \(V_m\).
\end{proposition}

\begin{proof}
Write \(R_\mu=z^\mu U\) for a finite-dimensional coefficient
space \(U\subset K_1\); positive finite generation ensures
finiteness. Choose rational \(m'>\|\beta\|\) with
\(m'\lesssim B'\), and a nonzero opposite section of weight
\(d'(m'\chi-\beta)\), where \(d'>0\) is an arbitrary
clearing factor as at the start of the twist proof.
Increase \(d'\) so that \(M=d'(m+m')\in n\Z_{>0}\).
If \(g'\) is its coefficient, every nonzero \(g\in U\) gives
\(z^{M\chi}g^{d'}g'\in R\). By \eqref{eq:A9},
\[
 D_1=H_1+\frac1M\Div_{S_1}
                    \bigl(g^{d'}g'/f_h^{M/n}\bigr)\ge0.
\]

Assume \(\dim S_1>0\). Construct compatible smooth projective
high models with morphisms down the tower, resolving rational
maps successively from the bottom, and include descents for
all displayed data. Choose compatible general closed points
where the morphisms are smooth. Lift relative parameters to
nested Taylor systems with \(p_i\) parameters at level \(i\).
Avoid the supports of the fixed forms, boundary, moduli and
\(H_1\). Make \(g',f_h\) units and a fixed basis of \(U\)
regular there. Give the parameters at level \(i\) rational
weights \(1/s_i\). This monomial test \(P_1\) is divisorial
up to scale: it is discrete up to scaling, and ratios of
equal-weight monomials give the required residue
transcendence degree. Its restriction \(P_i\) uses exactly
the parameters at levels \(\ell\ge i\).
At the chosen points the smooth monomial formula gives
\[
                  a_i(P_i)=s_i\sum_{\ell\ge i}p_\ell/s_\ell
                                      \lesssim1,
\]
using the bounded tower length and \(s_{\ell}\gtrsim s_i\).
Proposition~\ref{b:pure-upper} therefore implies
\[
                         P_1(g)\le CM/d'\le C'B'
\]
for every nonzero \(g\in U\).
Taylor truncation at this weighted bound is injective on \(U\);
the number of allowed monomials is at most
\(C\prod_i(1+B's_i)^{p_i}\).
If \(\dim S_1=0\), use \(K_1=k\). This proves \eqref{eq:C7}.

At fixed physical degree, \eqref{b:torus-axis-norm} bounds each
integer coordinate by \(|k_j|\le CS'/L_j\).
Counting these weights and using \eqref{eq:C7} proves
\eqref{eq:C8}. Weight zero contributes only constants.
The angular translation following \eqref{eq:A12} identifies
this count with the associated graded truncation of \(V_m\).
There are only finitely many costs below a fixed cutoff,
so it is also the dimension modulo higher costs.
\end{proof}

\begin{proposition}\label{b:actual-probes}
Each residual level has \(p_i\) actual probes in some \(V_m\),
of degree and cost \(O(1/s_i)\), whose initials are in \(K_i\)
and algebraically independent over \(K_{i+1}\).
Each torus axis has opposite angular probes in $V_{m_j}$ with initial
weights $\pm d_j\beta_j$ in the full initial field; their corresponding
degree-$m_j$ sections have weights $m_j\chi\pm d_j\beta_j$.
Here \(d_j\) are bounded positive integers,
\(m_j\in n\Z_{\ge0}\), and degrees and costs are \(O(1+L_j)\). The residual and positive torus initials together
are algebraically independent, and their monomials give box
lower bounds with widths \(\max(1,W_j)\).
\end{proposition}

\begin{proof}
The birational system in \eqref{eq:B3}, containing \(1\), gives
\(p_i\) functions algebraically independent over \(K_{i+1}\)
whose poles are bounded by \(kH_i\), \(k\lesssim1/s_i\).
For any one of them, \eqref{eq:B4} and pullback orders give
\[
             \Div_{S_1}(g)+kH_1\ge0.
\]
A trivial restriction evaluates both the lower function and
lower divisor at zero. Increase \(k\) to
\(m\in n\Z_{\ge0}\) with \(m\lesssim1/s_i\).
Then \eqref{eq:A9} gives \(g\bar h^{m/n}\in R_{m\chi}\).
Its angular lift lies in \(V_m\) with initial \(g\) and cost \(m\).

For a torus axis use \eqref{eq:C5} with
\(L=C\max(1,L_j)\), where the fixed uniform constant \(C\ge1\)
is chosen using \eqref{b:torus-axis-norm} so that \(\|\beta_j\|\le L\).
The resulting opposite sections lift to
actual angular functions after dividing their initials by
\(\bar h^{m_j/n}\). Their costs satisfy the asserted bound
because \(\widetilde v\in Q_T\).

Tower independence first makes the residual initials
algebraically independent. The positive torus initials then
have independent Laurent charges over \(K_1\), so adjoining
them preserves algebraic independence. For \(S'\ge1\),
take monomial exponents through
\[
                  \left\lfloor
                         \frac{cS'}{\max(1,W_j)}
                    \right\rfloor
\]
at each axis, for sufficiently small fixed \(c>0\).
Their degrees sum to at most \(S'\); nest the products into
a common allowed degree, rounding up to a multiple of \(n\).
Their costs are \(O(S')\), and their independent initials
ensure the same upper cost bound for nonzero combinations.
This proves the probe and box assertions.
\end{proof}
